% =====================================================================
%  Capítulo 9 — Violencia en la pareja: el agresor conocido
% =====================================================================
\chapter{Violencia en la pareja: el agresor conocido}

Un manual de defensa personal para mujeres quedaría incompleto sin abordar el escenario estadísticamente más probable: la violencia ejercida por la pareja o expareja. Aquí la \enquote{distancia preventiva} no es física sino temporal: detectar las señales cuanto antes.

\section{Señales de alarma tempranas}

\begin{itemize}
  \item Control: del móvil, la ropa, el dinero, las amistades, la ubicación.
  \item Aislamiento progresivo de familia y amigas.
  \item Celos presentados como \enquote{amor}; vigilancia de redes sociales.
  \item Humillaciones, desprecio a tus opiniones, cambios bruscos de humor.
  \item Empujones, golpes a objetos, amenazas\ldots{} seguidos de arrepentimiento y \enquote{luna de miel}. El ciclo de la violencia (tensión -- explosión -- arrepentimiento) tiende a acortarse y agravarse con el tiempo.
\end{itemize}

\section{Plan de seguridad}

\begin{itemize}
  \item Memoriza los teléfonos clave (\recurso{016}, \recurso{112}) y mantén el móvil cargado y contigo.
  \item Instala AlertCops y activa el Botón SOS (ver capítulo de aplicaciones): pulsándolo 5 veces en menos de 6 segundos envía alerta y audio a la policía.
  \item Prepara una bolsa de emergencia (documentación, llaves, dinero, medicación) en casa de alguien de confianza.
  \item Acuerda una palabra clave con familiares o vecinas que signifique \enquote{llama a la policía}.
  \item Identifica las zonas seguras de la casa (con salida, sin objetos peligrosos: evita cocina y baño durante las explosiones).
  \item Guarda pruebas: partes médicos, mensajes, fotografías de lesiones.
\end{itemize}

\section{Recursos específicos en España}

\begin{itemize}
  \item \recurso{016:} teléfono gratuito y confidencial, 24 horas, en varios idiomas. No deja rastro en la factura (recuerda borrarlo del registro de llamadas). También WhatsApp \recurso{600 000 016} y correo \href{mailto:016-online@igualdad.gob.es}{016-online@igualdad.gob.es}.
  \item \recurso{ATENPRO:} servicio telefónico de atención y protección con terminal móvil y localización, para víctimas con orden de protección (se solicita en servicios sociales).
  \item \recurso{Sistema VioGén:} seguimiento policial integral de los casos denunciados.
  \item \recurso{En emergencia:} \recurso{112}, \recurso{091} (Policía Nacional), \recurso{062} (Guardia Civil) o botón SOS de AlertCops.
\end{itemize}
